% 附录入口：代码清单，由完整论文主入口作为附录 C 引入。
\USTSRestoreEquationNumbers
\USTSAppendixChapter{代码}

% ！！！当代码文件含_时，请进行转义，如改为\_，以避免编译错误。！！！

\begin{lstlisting}[language=Matlab,caption={main.m/数据平滑与绘图代码},label={lst:appendix-matlab-main}]
clear; clc;

sampleCount = 180;
t = linspace(0, 12, sampleCount);
rng(2026);

signal = sin(t) + 0.35 * cos(2.4 * t);
noise = 0.18 * randn(size(t));
measurement = signal + noise;

windowSize = 9;
filtered = movmean(measurement, windowSize);
residual = measurement - filtered;
rmse = sqrt(mean((filtered - signal) .^ 2));

fprintf("RMSE = %.4f\n", rmse);

figure("Color", "w");
subplot(2, 1, 1);
plot(t, measurement, ".", "DisplayName", "measurement");
hold on;
plot(t, filtered, "LineWidth", 1.6, "DisplayName", "filtered");
plot(t, signal, "--", "LineWidth", 1.2, "DisplayName", "reference");
grid on;
xlabel("time/s");
ylabel("amplitude");
legend("Location", "best");

subplot(2, 1, 2);
plot(t, residual, "LineWidth", 1.2);
grid on;
xlabel("time/s");
ylabel("residual");
\end{lstlisting}

\begin{lstlisting}[language=Python,caption={preprocess.py/数据归一化代码},label={lst:appendix-normalize}]
from dataclasses import dataclass


@dataclass
class NormalizationStats:
    mean: float
    std: float


def normalize(values):
    if not values:
        raise ValueError("values must not be empty")
    mean = sum(values) / len(values)
    variance = sum((item - mean) ** 2 for item in values) / len(values)
    std = variance ** 0.5 or 1.0
    return [(item - mean) / std for item in values], NormalizationStats(mean, std)


scores, stats = normalize([71.5, 83.0, 79.5, 92.0])
print(scores)
print(stats)
\end{lstlisting}

\begin{lstlisting}[language=C++,caption={search.cpp/网格搜索代码},label={lst:appendix-grid}]
#include <iostream>
#include <limits>
#include <vector>

double objective(double x, double y) {
    return (x - 3.0) * (x - 3.0) + (y + 2.0) * (y + 2.0);
}

int main() {
    double best_value = std::numeric_limits<double>::infinity();
    double best_x = 0.0;
    double best_y = 0.0;

    for (double x = -5.0; x <= 5.0; x += 0.5) {
        for (double y = -5.0; y <= 5.0; y += 0.5) {
            double value = objective(x, y);
            if (value < best_value) {
                best_value = value;
                best_x = x;
                best_y = y;
            }
        }
    }

    std::cout << "best: " << best_x << ", " << best_y << std::endl;
    return 0;
}
\end{lstlisting}

\begin{lstlisting}[language=SQL,caption={report.sql/实验结果汇总代码},label={lst:appendix-sql}]
SELECT
    model_name,
    COUNT(*) AS run_count,
    ROUND(AVG(accuracy), 4) AS avg_accuracy,
    ROUND(AVG(f1_score), 4) AS avg_f1
FROM experiment_runs
WHERE finished_at >= DATE '2026-01-01'
GROUP BY model_name
HAVING COUNT(*) >= 3
ORDER BY avg_f1 DESC, avg_accuracy DESC;
\end{lstlisting}
